% !PRINTED 466
\weberpage{466}

%<<CONT
durch $\Omega(\varrho)$ bezeichnen. Ist $\mu$ der Grad von $N$, so genügt $\varrho$ einer irreducibeln Gleichung $\mu^{\mathrm{ten}}$ Grades
\weq{5}{g(t) = 0,}
von der zu zeigen ist, dass es eine Normalgleichung ist. Zu diesem Zweck bemerken wir zunächst, dass die eine Wurzel $\varrho$ dieser Gleichung eine rationale Function der $\alpha, \alpha_1 \ldots \alpha_{m-1}$ ist, weil sie in $N$ enthalten war. Setzen wir, um dies anzudeuten,
\[
\varrho = \varrho(\alpha, \alpha_1 \ldots \alpha_{m-1}),
\]
und bilden nun alle verschiedenen Anordnungen der Ziffern
\[
0,\ 1,\ 2 \ldots m - 1,
\]
deren Anzahl $\Pi(m)$ beträgt:
\weq{6}{(0, 1, 2 \ldots m - 1),\ (0', 1', 2' \ldots (m-1)'),\ (0'', 1'', 2'' \ldots (m-1)'')\,\ldots,}
worin die Ziffern mit einem, zwei etc. Accenten dieselben sind, wie die ohne Accent, nur in anderer Reihenfolge, und bilden hieraus die Functionen
\weq{7}{\varrho = \varrho(\alpha_0, \alpha_1 \ldots \alpha_{m-1}),\ \varrho' = \varrho(\alpha_{0'}, \alpha_{1'} \ldots \alpha_{(m-1)'}),\ \varrho'' = \varrho(\alpha_{0''}, \alpha_{1''} \ldots \alpha_{(m-1)''})\,\ldots,}
unbekümmert darum, ob darunter etwa unter einander gleiche vorkommen oder nicht.

Wenn wir in allen den Anordnungen (6) ein und dieselbe Vertauschung vornehmen, z.\,B. 0 mit 1, so ändert sich die Gesammtheit dieser Anordnungen nicht, sondern nur ihre Reihenfolge wird eine andere. Denn erstens kann durch eine solche Vertauschung nichts Anderes entstehen, als Anordnungen der Ziffern, und zweitens können nicht zwei verschiedene Anordnungen durch eine und dieselbe Vertauschung in dieselbe Anordnung übergehen.

Es werden also auch durch jede solche Vertauschung die Functionen (7) nur unter einander permutirt werden. Bilden wir also das Product
\[
G(t) = (t - \varrho)\,(t - \varrho')\,(t - \varrho'')\,\ldots
\]
für eine Veränderliche $t$, so bleiben seine Coëfficienten, die gewiss Functionen von $\alpha, \alpha_1 \ldots \alpha_{m-1}$ sind, ungeändert, wenn diese Grössen irgendwie permutirt werden; d.\,h. es sind symmetrische Functionen der Wurzeln der Function $F(x)$, und also ist $G(t)$ eine Function von $t$ in $\Omega$. Alle Wurzeln von $G(t)$ sind Grössen in $N$, da sie durch die $\alpha$ rational ausgedrückt sind.
